\chapter{Backpropagation and Automatic Differentiation}

\section{Computational Graphs}

Neural networks are compositions of functions. To compute gradients efficiently, it is useful to represent these computations explicitly.

\medskip

\noindent
\textbf{Computational graph.}  
A computational graph is a directed acyclic graph (DAG) that represents how a function is built from simpler operations.

\medskip

\noindent
\textbf{Nodes and edges.}
\begin{itemize}
    \item Nodes represent variables or intermediate values
    \item Edges represent functional dependencies
\end{itemize}

\medskip

\noindent
\textbf{Example.}  
Consider a simple function
\[
z = (x_1 + x_2) \cdot x_3.
\]

\medskip

\noindent
We can represent this as a graph:
\begin{itemize}
    \item $a = x_1 + x_2$
    \item $z = a \cdot x_3$
\end{itemize}

\medskip

\noindent
\textbf{Forward pass.}  
We evaluate the graph from inputs to outputs, computing all intermediate values.

\medskip

\noindent
\textbf{Insight.}  
A neural network is a large computational graph, where each layer corresponds to a set of operations.

\section{Chain Rule and Reverse-Mode Differentiation}

To train a neural network, we must compute gradients of a loss function with respect to all parameters.

\medskip

\noindent
\textbf{Chain rule.}  
For a composition of functions,
\[
f(x) = f_L(f_{L-1}(\cdots f_1(x))),
\]
the derivative is obtained by repeated application of the chain rule.

\medskip

\noindent
\textbf{Local derivatives.}  
Each node in the computational graph contributes a local derivative.

\medskip

\noindent
\textbf{Reverse-mode differentiation.}  
Gradients are computed by propagating sensitivities from the output backward through the graph.

\medskip

\noindent
\textbf{Basic principle.}  
If a variable $z$ depends on $x$ through intermediate variables, then
\[
\frac{\partial L}{\partial x} = \sum_{\text{paths}} \frac{\partial L}{\partial z} \cdot \frac{\partial z}{\partial x}.
\]

\medskip

\noindent
\textbf{Backpropagation.}  
This process of propagating gradients backward through the network is known as backpropagation.

\medskip

\noindent
\textbf{Example.}  
For $z = a \cdot x_3$ and $a = x_1 + x_2$,
\[
\frac{\partial z}{\partial x_1} = \frac{\partial z}{\partial a} \cdot \frac{\partial a}{\partial x_1}.
\]

\medskip

\noindent
\textbf{Insight.}  
Backpropagation systematically applies the chain rule across the computational graph.

\section{Efficient Gradient Computation}

A naive computation of gradients would be computationally expensive, especially for deep networks.

\medskip

\noindent
\textbf{Key idea.}  
Reuse intermediate computations from the forward pass.

\medskip

\noindent
\textbf{Forward pass.}
\begin{itemize}
    \item Compute all intermediate values
    \item Store them for reuse
\end{itemize}

\medskip

\noindent
\textbf{Backward pass.}
\begin{itemize}
    \item Compute gradients in reverse order
    \item Reuse stored intermediate values
\end{itemize}

\medskip

\noindent
\textbf{Complexity.}  
Backpropagation computes gradients with a cost comparable to evaluating the function itself.

\medskip

\noindent
\textbf{Gradient flow.}  
Gradients propagate layer by layer:
\[
\frac{\partial L}{\partial W_\ell}
\quad \text{depends on} \quad
\frac{\partial L}{\partial z^{(\ell)}}.
\]

\medskip

\noindent
\textbf{Shared subcomputations.}  
The computational graph ensures that repeated operations are computed only once.

\medskip

\noindent
\textbf{Insight.}  
Backpropagation is efficient because it avoids redundant computation and exploits the structure of the network.

\section{Implementation Perspectives}

Modern machine learning frameworks implement backpropagation using automatic differentiation.

\medskip

\noindent
\textbf{Automatic differentiation (AD).}
\begin{itemize}
    \item Computes exact derivatives (up to numerical precision)
    \item Works for arbitrary compositions of differentiable operations
\end{itemize}

\medskip

\noindent
\textbf{Two modes of AD.}
\begin{itemize}
    \item Forward mode
    \item Reverse mode (used in deep learning)
\end{itemize}

\medskip

\noindent
\textbf{Why reverse mode?}  
Efficient when the function has many parameters and a single output (typical in learning problems).

\medskip

\noindent
\textbf{Practical workflow.}
\begin{enumerate}
    \item Define the model as a computational graph
    \item Perform a forward pass to compute loss
    \item Use automatic differentiation to compute gradients
    \item Update parameters using an optimization algorithm
\end{enumerate}

\medskip

\noindent
\textbf{Memory considerations.}
\begin{itemize}
    \item Intermediate values must be stored for backward pass
    \item Large models require careful memory management
\end{itemize}

\medskip

\noindent
\textbf{Numerical issues.}
\begin{itemize}
    \item Vanishing and exploding gradients
    \item Finite precision arithmetic
\end{itemize}

\medskip

\noindent
\textbf{Insight.}  
Automatic differentiation abstracts the complexity of gradient computation, allowing practitioners to focus on model design.

\section{A Unifying Perspective}

Backpropagation connects several core ideas:

\begin{itemize}
    \item \textbf{Mathematics:} chain rule for derivatives
    \item \textbf{Computation:} efficient reuse of intermediate values
    \item \textbf{Algorithms:} scalable training of deep networks
\end{itemize}

\medskip

\noindent
\textbf{Key components.}
\begin{itemize}
    \item Computational graph
    \item Forward pass
    \item Backward pass
\end{itemize}

\medskip

\noindent
\textbf{Core insight.}  
Backpropagation enables efficient training of complex models by systematically applying the chain rule.

\section{Outlook}

This chapter introduced backpropagation and automatic differentiation:
\begin{itemize}
    \item computational graphs as a representation of models,
    \item reverse-mode differentiation for efficient gradients,
    \item practical implementation via automatic differentiation.
\end{itemize}

\medskip

In the next chapter, we study the challenges of training deep networks, including initialization and gradient-related issues.

\medskip

\noindent
\textbf{Guiding principle.}  
Efficient computation of gradients is the key that makes modern deep learning possible.